\documentclass[11pt,a4paper]{article}

\usepackage[utf8]{inputenc}
\usepackage[T1]{fontenc}
\usepackage[margin=2.5cm]{geometry}
\usepackage{parskip}
\usepackage{amsmath,amssymb}
\usepackage{enumitem}
\usepackage{xcolor}
\usepackage{hyperref}
\hypersetup{colorlinks=true,linkcolor=blue,urlcolor=blue}

\newenvironment{quotebox}
  {\begin{quote}\itshape\color{black!70}}
  {\end{quote}}

\title{Should I give up publishing the cut-and-project paper?\\[4pt]
       \large A strategic assessment}
\author{Notes prepared with Claude Code}
\date{2026-06-09}

\begin{document}
\maketitle

\noindent
\textit{Context: after sending the paper and a prior-art note to Dr.\ Henna
Koivusalo (Bristol), she kindly declined --- agreeing with Alan Haynes that
the rational cut-and-project case is ``considered solved'' / folklore, while
noting that the Lean formalization may be new. The question: should I give up
on publishing? What follows is an honest assessment, not encouragement.}

\section*{My honest read}

You shouldn't give up. But you should give up on \textbf{one specific thing},
and that clarity is what unlocks the rest.

\textbf{Give up on:} \emph{``the theorem is a novel number-theory result,
recognized as such by number theorists.''} Two independent experts in exactly
this subfield (Haynes and Koivusalo, who is part of the
Haynes--Koivusalo--Walton--Sadun group whose paper is your closest pendant)
have now told you the same thing, consistently, and Koivusalo even handed you
the references where the folklore lives. Continuing to pitch the \emph{theorem}
as new to that circle will cost you time and credibility, and it won't change
their verdict. That door is honestly closed.

\textbf{Do not give up on:} \emph{getting this work published, cited, and
recognized.} That's a different door, and Koivusalo's email actually props it
open rather than shutting it. Read her second-to-last paragraph again:

\begin{quotebox}
I do also agree that the lean application may be new.
\end{quotebox}

That is the same thing Haynes said unprompted. Two domain experts, neither of
whom you asked about Lean, both volunteered that the formalization is the part
with novelty. That's not a consolation prize --- that's them pointing at your
actual contribution.

\section*{So the pivot isn't a retreat --- it's what they're telling you to do}

Your own strategy notes (\texttt{support/lean.md}) already nailed the reframe a
week ago:

\begin{quotebox}
\textbf{Not:} ``I proved a new theorem, please recognize it.''\\
\textbf{Instead:} ``I formalized X --- is it useful, is it already in Mathlib,
what did I miss?''
\end{quotebox}

Koivusalo's reply is evidence \emph{for} that pivot, not against it. The paper
doesn't die; it changes what it claims to be. The honest framing of the paper
becomes:

\begin{quotebox}
A fully machine-checked (Lean 4, 0 \texttt{sorry}/\texttt{axiom}) treatment of
the gap-period of 1-D rational cut-and-project sequences. The underlying
number-theoretic result is classical/folklore; the contribution is the
formalization and the reusable residue-distribution lemmas.
\end{quotebox}

That is a paper that is \textbf{honest, defensible, and publishable} --- just
not in the venue you originally aimed at.

\section*{What ``publish'' realistically means now, in priority order}

These are straight from your own notes, and Koivusalo just validated the
premise behind all of them:

\begin{enumerate}[leftmargin=*]
  \item \textbf{arXiv preprint, reframed around the formalization.} Legitimate
  regardless of theorem-novelty. The one change: seek an endorser from
  \textbf{cs.LO / math.LO}, not number theory. A formalization framing fits
  that category far better anyway.

  \item \textbf{Mathlib contribution (your Tier A residue-distribution
  lemmas).} This is the most concrete \emph{recognition} path you have ---
  merit-based, no endorsement gatekeeping, just review. Your Mathlib pre-check
  suggested those lemmas aren't already there. If they land, you have a
  permanent, cited contribution to a tool thousands of people use. That's a
  more durable form of ``published'' than a journal article almost no one
  reads.

  \item \textbf{Blog post + Lean Zulip intro.} Cheapest, highest-leverage way
  to make the artifact visible.

  \item \textbf{CPP / ITP short paper} --- possible later, but be honest with
  yourself: a 1-D folklore result is borderline as a \emph{standalone}
  formalization paper. Don't anchor on it.
\end{enumerate}

\section*{Two pieces of honest homework before you publish anything}

\begin{enumerate}[leftmargin=*]
  \item \textbf{Actually verify Koivusalo's two references.} Open Haynes's
  \emph{Aperiodic Order} lecture notes and Pytheas Fogg (\emph{Substitutions in
  Dynamics, Arithmetics and Combinatorics}, Springer LNM 1794) and check
  whether the period formula is genuinely stated there. This matters either
  way: if it's there, you cite it and your prior-art section becomes
  bulletproof; if it's \emph{not} explicitly there, that actually
  \emph{strengthens} your ``we wrote down and formalized something only
  folklore-known'' framing. You can't write the honest version of the paper
  without doing this check.

  \item \textbf{Pin down the set-valued case (Open Problem 2).} Your literature
  search couldn't find an explicit closed formula for the minimal period of the
  \emph{set-valued} (multiplicity-discarded) gap sequence --- the three-gap
  theorem gives structure but not the minimal period. This is the one sliver
  where I'd genuinely keep an open mind about novelty. But be disciplined: the
  experts said ``the rational case is solved,'' and they may well consider this
  part folklore too. Don't let it become the thing you cling to. Treat it as a
  precise question worth one careful check, not a rescue.
\end{enumerate}

\section*{Bottom line}

Your instinct not to give up is right --- and you don't have to. But the version
of ``not giving up'' that works is the one where you \emph{accept the experts'
verdict on the theorem} and \emph{act on the verdict they volunteered about the
Lean work}. That's not surrender; it's aiming at the target they're both
pointing at. The cruel version of this would be to spend three more months
trying to get number theorists to call the formula new. The kind version ---
to yourself and the work --- is to put the formalization in front of the people
who can actually judge it and use it.

\end{document}
